\documentclass[12pt, a4paper]{article}
\usepackage[english,spanish,es-tabla ]{babel}
\usepackage[top=0.5cm, left=0.5cm, right=0.5cm, bottom=0.5cm]{geometry}
\usepackage{mystyle_2}
\usepackage{mycolours}
\usepackage{mytcolorbox}
\usepackage{mycommands}
\usepackage{microtype}
\usepackage[all]{chemmacros}

\usepackage{pdflscape}
\usepackage{rotating}
\usepackage{tabularray}
\UseTblrLibrary{siunitx}
\newcommand{\acu}{{(aq)}}

\begin{document}

\DefTblrTemplate{contfoot-text}{normal}{Continúa en la siguiente página}
\SetTblrTemplate{contfoot-text}{normal}
\DefTblrTemplate{conthead-text}{normal}{(Continuación)}
\SetTblrTemplate{conthead-text}{normal}

\begin{landscape}

    \begin{longtblr}[caption=Oxoácidos]{hlines,vlines, colspec={Q[c,m] Q[c,m, 2cm] Q[c,m,7cm] Q[c,m,7cm]}, rowsep=5pt, rowhead=1}
        \SetRow{gray!50} nº  & Fórmula  & Tradicional  & Hidrógeno  \\
        1 & \ch{HBrO}  &   & \\
        2 & \ch{HBrO2}  &   & \\
        3 & \ch{HBrO3}  &   & \\
        4 & \ch{HBrO4}  &   & \\
        5 & \ch{HIO}  &   & \\
        6 & \ch{HIO2}  &   & \\
        7 & \ch{HIO3}  &   & \\
        8 & \ch{HIO4}  &   & \\
        9 & \ch{HNO2}  &   & \\
        10 & \ch{H2SeO2}  &   & \\
        11 & \ch{H2SO3}  &   & \\
        12 & \ch{H2TeO2}  &   & \\
        13 & \ch{H2SO4}  &   & \\
        14 & \ch{H2TeO3}  &   & \\
        15 & \ch{HClO}  &   & \\
        16 & \ch{HClO2}  &   & \\
        17 & \ch{HClO3}  &   & \\
        18 & \ch{HClO4}  &   & \\
        19 & \ch{HNO3}  &   & \\
        20 & \ch{HNO2}  &   & \\
        21 & \ch{H2SO2}  &   & \\
        22 & \ch{H2SO3}  &   & \\
        23 & \ch{H2SO4}  &   & \\
        24 & \ch{HMnO4}  &   & \\
        25 & \ch{H3PO3}  &   & \\
        26 & \ch{H2CrO3}  &   & \\
        27 & \ch{H3PO4}  &   & \\
        28 & \ch{H3AsO4}  &   & \\
        29 & \ch{HSbO2}  &   & \\
        30 & \ch{H3SbO3}  &   & \\
        31 & \ch{HSbO3}  &   & \\
        32 & \ch{H3SbO4}  &   & \\
        33 & \ch{HNO}  &   & \\
        34 & \ch{H2CO3}  &   & \\
        35 & \ch{H2CO2}  &   & \\
        36 & \ch{H2SeO4}  &   & \\
        37 & \ch{HMnO4}  &   & \\
        38 & \ch{H2MnO4}  &   & \\
        39 & \ch{H2MnO3}  &   & \\
        40 & \color{red} \ch{}  & ácido hipocloroso  &\color{red} \\
        41 & \color{red} \ch{}  & ácido hipyodoso  &\color{red} \\
        42 & \color{red} \ch{}  & ácido nítrico  &\color{red} \\
        43 & \color{red} \ch{}  &  ácido selénico &\color{red} \\
        44 & \color{red} \ch{}  &  ácido sulfúrico &\color{red} \\
        45 & \color{red} \ch{}  &  ácido telúrico &\color{red} \\
        46 & \color{red} \ch{}  & ácido cloroso  &\color{red} \\
        47 & \color{red} \ch{}  & ácido perclórico  &\color{red} \\
        48 & \color{red} \ch{}  & ácido hipocloroso  &\color{red} \\
        49 & \color{red} \ch{}  &  ácido clórico &\color{red} \\
        50 & \color{red} \ch{}  &  ácido hipobromoso  &\color{red} \\
        51 & \color{red} \ch{}  &  ácido bromoso  &\color{red} \\
        52 & \color{red} \ch{}  &  ácido brómico &\color{red} \\
        53 & \color{red} \ch{}  & ácido perbrómico   &\color{red} \\
        54 & \color{red} \ch{}  &  ácido mangánico  &\color{red} \\
        55 & \color{red} \ch{}  &  ácido permangánico  &\color{red} \\
        56 & \color{red} \ch{}  &  ácido crómico  &\color{red} \\
        57 & \color{red} \ch{}  &  ácido hiposulfuroso  &\color{red} \\
        58 & \color{red} \ch{}  &  ácido sulfúrico   &\color{red} \\
        59 & \color{red} \ch{}  &  ácido telúrico  &\color{red} \\
        60 & \color{red} \ch{}  & ácido nitroso   &\color{red} \\
        61 & \color{red} \ch{}  &  ácido fosfórico  &\color{red} \\
        62 & \color{red} \ch{}  &   ácido nítrico &\color{red} \\
        63 & \color{red} \ch{}  &  ácido sulfuroso  &\color{red} \\
        64 & \color{red} \ch{}  &  ácido teluroso  &\color{red} \\
        65 & \color{red} \ch{}  &  ácido selénico  &\color{red} \\
        66 & \color{red} \ch{}  &  ácido carbónico  &\color{red} \\
        67 & \color{red} \ch{}  &  ácido carbonoso  &\color{red} \\
        68 & \color{red} \ch{}  &  ácido silícico  &\color{red} \\
        69 & \color{red} \ch{}  &   ácido arsenioso &\color{red} \\
        70 & \color{red} \ch{}  &  \color{red}  & Hidrogeno(dioxidobromato)\\
        71 & \color{red} \ch{}  & \color{red}   & Hidrogeno(oxidobromato) \\
        72 & \color{red} \ch{}  &  \color{red}  &  Hidrogeno(trioxidobromato) \\
        73 & \color{red} \ch{}  &  \color{red}  & Hidrogeno(tetraoxidobromato) \\
        74 & \color{red} \ch{}  &  \color{red}  & Hidrogeno(oxidoyodato)\\
        75 & \color{red} \ch{}  &  \color{red}  & Hidrogeno(dioxidoyoadato)\\
        76 & \color{red} \ch{}  &  \color{red}  &  Hidrogeno(trioxidoyodato)\\
        77 & \color{red} \ch{}  &  \color{red}  & Hidrogeno(tetraoxidoyodato)\\
        78 & \color{red} \ch{}  &  \color{red}  & Hidrogeno(dioxidonitrato)\\
        79 & \color{red} \ch{}  & \color{red}   &Dihidrogeno(dioxidoseleniato) \\
        80 & \color{red} \ch{}  &  \color{red}  &Dihidrogeno(trioxidosulfato)\\
        81 & \color{red} \ch{}  &  \color{red}  &Dihidrogeno(tetraoxidosulfato) \\
        82 & \color{red} \ch{}  &  \color{red}  &Dihidrogeno(trioxidotelurato) \\
        83 & \color{red} \ch{}  &  \color{red}  &Dihidrogeno(tetraoxidotelurato) \\
        84 & \color{red} \ch{}  &  \color{red}  &Hidrogeno(oxidoclorato) \\
        85 & \color{red} \ch{}  &  \color{red}  &Hidrogeno(dioxidoclorato) \\
        86 & \color{red} \ch{}  &  \color{red}  &Hidrogeno(trioxidoclorato) \\
        87 & \color{red} \ch{}  &  \color{red}  &Hidrogeno(tetraoxidobromato) \\
        88 & \color{red} \ch{}  &  \color{red}  &Hidrogeno(oxidobromato) \\
        89 & \color{red} \ch{}  &  \color{red}  &Hidrogeno(dioxidobromato) \\
        90 & \color{red} \ch{}  &  \color{red}  &Dihidrogeno(tetraoxidoseleniato) \\
        91 & \color{red} \ch{}  &  \color{red}  &Dihidrogeno(tetraoxidosulfato)\\
        92 & \color{red} \ch{}  &  \color{red}  &Dihidrogeno(tetraoxidotelurato) \\
    \end{longtblr}

\end{landscape}
\end{document}